\resumeAn
%
As part of my Master's degree in Networks and Computer Systems, I completed my graduation project at \textbf{Inetum Offshore Maroc}, within the Telecom Business Unit. The project focuses on the design and implementation of an \textbf{AI agent for technical diagnosis} targeting the \textbf{CCU} (Commande et Catalogue Unifiés) information system of a telecom operator.

\medskip

The objective is to accelerate and improve the reliability of technical incident handling by automatically correlating multiple data sources: \textbf{network logs}, \textbf{customer context} from the CRM, and \textbf{historical tickets} stored in a \textbf{Neo4j} knowledge graph. The system relies on an agentic orchestration built with \textbf{LangGraph}, an extensible embedding interface, and a local language model via \textbf{Ollama}.

\medskip

The first part of the work focused on designing the agentic architecture: incident parsing, parallel source collection, root-cause reasoning, remediation proposal, content guardrail, PDF report generation, and notification. The second part covered the technical implementation: Python implementation with \textbf{FastAPI}, \textbf{Neo4j} graph database including a vector index, ingestion of mocked data, \textbf{pytest} testing, and containerization via \textbf{Docker Compose}.

\medskip

This project allowed me to address the challenges of generative AI applied to telecom support, multi-agent system design, graph-augmented retrieval (\textbf{GraphRAG}), and enterprise LLM deployment, while respecting a strong constraint: no corrective action is ever executed automatically on a real system.


\vspace{1cm}


\noindent\rule[2pt]{\textwidth}{0.5pt}

{\textbf{Keywords:}}
AI Agent, Technical Diagnosis, CCU, LangGraph, Neo4j, GraphRAG, FastAPI, TMF620, TMF622, LLM, Ollama, Telecom Support.
\\
\noindent\rule[2pt]{\textwidth}{0.5pt}
